\documentclass[11pt]{article}
\usepackage[utf8]{inputenc}
\usepackage[T1]{fontenc}
\usepackage{amsmath,amssymb}
\usepackage{booktabs}
\usepackage{graphicx}
\usepackage{microtype}
\usepackage[margin=1in]{geometry}
\usepackage[numbers,sort&compress]{natbib}
\usepackage{xcolor}
\usepackage[colorlinks=true,allcolors=blue]{hyperref}
\usepackage{multirow}

\newcommand{\dauroc}{$\Delta$AUROC}
\newcommand{\sae}{SAE}

\title{\bfseries SAE Features Are Causally Active but Predictively Redundant
for Difficulty:\\ Evidence Across Language and Time-Series Foundation Models}

\author{%
  Nabin Prasad Dev \\
  \texttt{github.com/nabindev3/fm-difficulty-probe}
}
\date{}

\begin{document}
\maketitle

\begin{abstract}
Foundation models ship without a native abstention signal, yet some inputs are
far harder than others. We ask whether a model's internal TopK sparse-autoencoder
(\sae) features encode a \emph{self-difficulty} signal beyond a cheap, no-forward-pass
baseline, and we ask it of two unrelated modalities through a single shared
pipeline: an autoregressive language model (Pythia-410M on HellaSwag and SQuAD)
and an encoder-based time-series foundation model (Chronos-T5-small on ETTh1).
Across both modalities, \sae{} features add \emph{no} incremental predictive power
for difficulty over the strongest cheap baseline (\dauroc{}(\sae{}$-$raw)
$\le 0$; label-permutation $p \le 0.03$, and $p<10^{-4}$ on SQuAD), and the
result is invariant to the \sae{} expansion factor (4$\times$ vs.\ 8$\times$).
The deployable consequence also replicates: a Platt-recalibrated selective
predictor built on the cheap baseline captures \mbox{30--41\%} of the oracle's
risk--coverage improvement and yields a Pareto-dominating cost--quality cascade
in both modalities. The causal picture, however, \emph{dissociates}: under
reconstruction-patching ablation the language model's top features are causally
active (5/5 significant under all-position intervention, 0--2/5 under
single-position---coverage, not fidelity), whereas the time-series model's
features are causally silent under either intervention coverage (0/5). The
cross-modal dissociation---predictive null in both, causal positive in the LM
only---localizes the causal difficulty signal as a property of the autoregressive
language model rather than a universal property of sparse features. We release a
modality-agnostic, regression-gated, one-command-reproducible implementation.
\end{abstract}

\section{Introduction}
\label{sec:intro}

Foundation models (FMs) are deployed as black boxes: a language model (LM)
generates an answer, a time-series foundation model (TSFM) emits a forecast, and
in both cases every input consumes the same compute and is granted the same
trust, even though some inputs are several times harder than others. Knowing
\emph{which} inputs a model will struggle on---before paying for an expensive
model or committing to a forecast---is the core of routing, abstention, and
cost--quality cascades \citep{chen2023frugalgpt,ong2024routellm}.

A natural hypothesis, given the rise of mechanistic interpretability, is that an
FM's internal representations already encode this self-difficulty signal, and
that sparse autoencoders (\sae{}s)---which decompose activations into
interpretable, monosemantic features \citep{bricken2023monosemanticity,gao2024scaling}---expose
it in a usable form. The open, and to our knowledge unclaimed, question is
whether \sae{} features add difficulty signal \emph{beyond what a cheap baseline
already captures}: lexical statistics for text, classical descriptive statistics
for a time series. A single negative answer on one model invites the reflex
``you just probed it wrong.'' We therefore ask the question of two modalities
\emph{through the same pipeline, with the same metrics}, so that what replicates
is a property of \sae{}s and what does not is localized to one architecture.

\paragraph{Contributions.}
(i) A modality-agnostic pipeline---one tested core, two thin adapters---that runs
the identical probe ladder, causal ablation, and selective-prediction analysis on
an LM and a TSFM, gated on reproducing both prior studies' headline numbers
before any new experiment is trusted (Section~\ref{sec:method}).
(ii) A \emph{predictive null} that replicates across modalities and is robust to
\sae{} width and to a label-permutation test (Sections~\ref{sec:llm}--\ref{sec:synth}).
(iii) A \emph{causal dissociation}: the same reconstruction-patching ablation
finds the LM's features causally active but the TSFM's features silent, and on
the LM side localizes per-feature detectability to intervention \emph{coverage},
not \sae{} fidelity (Section~\ref{sec:synth}).
(iv) A deployable artifact---a Platt-recalibrated selective predictor on the cheap
baseline---that replicates in both modalities (Section~\ref{sec:deploy}).

\section{Method: one pipeline, two modalities}
\label{sec:method}

We describe the pipeline once; the two backends differ only in an adapter that
implements a common \texttt{Modality} interface (cheap features, raw activations,
\sae{} codes, labels, splits, CV scheme, causal metric, cascade models). Nothing
in the shared core imports a specific model.

\paragraph{Backbones and data.}
The LM modality uses Pythia-410M \citep{biderman2023pythia} (residual stream at
Layer~12 ``mid'' and Layer~18 ``late''), evaluated on HellaSwag
\citep{zellers2019hellaswag} (binary multiple-choice correctness; $N{=}5000$,
3500/1500 train/test) and SQuAD \citep{rajpurkar2016squad} (continuous
gold-answer perplexity difficulty). The TSFM modality uses Chronos-T5-small
\citep{ansari2024chronos} (encoder block~3 ``mid'' and block~5 ``late'')
on ETTh1 \citep{zhou2021informer} ($701$ forecast windows, 483/167 train/test,
context $512$, horizon $96$). Cascades use Pythia-2.8B and Chronos-T5-base as the
expensive model.

\paragraph{Sparse autoencoder.}
A TopK \sae{} with $k{=}32$ and auxiliary dead-feature revival is trained on the
\emph{train split only}. The LM \sae{} uses $4\times$ expansion ($d_{\text{hidden}}{=}4096$),
the TSFM \sae{} $8\times$ ($d_{\text{hidden}}{=}4096$); we test invariance to this
choice in Section~\ref{sec:synth}.

\paragraph{The probe ladder.}
The crux of an apples-to-apples cross-modal claim is that both modalities report
the \emph{same} comparison. We fit an $L_1$-regularized logistic probe on five
feature sets and rank them: \textbf{P1} cheap baseline (lexical stats $\mid$
classical TS stats); \textbf{P2} cheap $+$ raw activations; \textbf{P3} cheap $+$
\sae{} codes; and diagnostic isolations \textbf{P4} raw-only, \textbf{P5}
\sae{}-only. The headline quantity is $\Delta$(\sae{}$-$raw) $=$ AUROC(P3)$-$AUROC(P2):
does the \sae{} decomposition add anything \emph{on top of} the raw activations it
was trained on? Regularization $C$ is chosen by inner cross-validation that
respects each modality's structure (stratified-by-category for the LM,
\texttt{TimeSeriesSplit} for the TSFM). All deltas carry paired-bootstrap 95\%
CIs ($B{=}2000$), and the headline delta additionally carries a label-permutation
$p$-value ($B{=}10{,}000$).

\paragraph{Causal ablation.}
We rank the top-5 difficulty-predictive latents and, for each, measure the change
in a continuous metric (negative log-probability / cross-entropy for the LM in
nats; CRPS \citep{gneiting2007strictly} for the TSFM) when that latent is zeroed
in a forward-hook \sae{} reconstruction of the chosen layer, relative to the
reconstruction baseline. Crucially, the intervention \emph{coverage} is a
first-class knob: \emph{single-position} (one boundary token / last context step)
versus \emph{all-position} (every prompt token / every context step), with the
metric held fixed so coverage is not confounded with the binary/continuous choice.

\paragraph{Selective prediction, calibration, cascade.}
Risk--coverage curves report the area under the curve (AURC) and the fraction of
the oracle's improvement captured, on each modality's natural error scale (binary
correctness for the LM, continuous CRPS for the TSFM). Probe probabilities are
recalibrated with Platt \citep{platt1999probabilistic} and isotonic regression
fit on 5-fold out-of-fold train predictions. A threshold sweep routes between the
cheap and expensive model, tracing a Pareto frontier against random and oracle
routing.

\paragraph{Leakage controls and the regression gate.}
Both adapters enforce: train-only \sae{} fitting, a temporal purge gap
$\ge$ context$+$horizon (TSFM) and pretraining-contamination / lexical-dedup
filters (LM), a prompt-only perplexity feature that never sees the gold target,
and CV that respects category/temporal structure. We adopt a strict
\emph{regression gate}: no new experiment is trusted until both prior studies'
headline numbers reproduce through the shared code---LM SQuAD/L18 raw-only
AUROC $=0.716$ and TSFM $\Delta$(\sae{}$-$cheap) $=-0.228$, both reproduced
exactly. The probe refuses to run on a missing or randomly-initialized \sae{}
checkpoint and on single-class labels.

\section{Results A --- Language model (Pythia)}
\label{sec:llm}

\paragraph{Predictive null.}
On SQuAD, raw residual activations are genuinely useful---P2 reaches AUROC $0.671$
at Layer~12 and $0.716$ at Layer~18, versus $0.590$ for lexical stats---but the
\sae{} adds nothing on top: $\Delta$(\sae{}$-$raw) $=-0.079$ (95\% CI
$[-0.118,-0.041]$), with a label-permutation $p<10^{-4}$. HellaSwag is a
near-chance regime for every probe family (P1 $0.509$, P3 $0.500$), so no method
adds signal there. In neither case does the \sae{} decomposition beat the raw
activations.

\paragraph{Causal positive, and coverage-not-fidelity.}
The same features the probe ranks are causally active: under all-position
patching, 5/5 of the top SQuAD features (and 5/5 HellaSwag) show ablation effects
whose CIs exclude zero. Under single-position patching the count collapses to
0--2/5---the effect is real but \emph{below the resolution of a single-token
intervention}. Per-feature causal attribution requires multi-position coverage,
not a higher-fidelity \sae{}.

\section{Results B --- Time-series model (Chronos-T5)}
\label{sec:tsfm}

\paragraph{Predictive null, more starkly.}
On ETTh1 the cheap classical statistics are the \emph{strongest} probe
(P1 $0.654$); raw encoder activations do worse (P2 $0.584$) and the \sae{} worse
still (P3 $0.426$). Thus $\Delta$(\sae{}$-$cheap) $=-0.228$ (CI $[-0.366,-0.092]$,
the legacy headline) and $\Delta$(\sae{}$-$raw) $=-0.158$ (CI $[-0.293,-0.025]$,
$p=0.011$). The null holds at both encoder layers.

\paragraph{Causal silence.}
Under the same reconstruction-patching protocol at full sample budget
($50$ samples/window, all $167$ test windows), \emph{no} feature is significant
under either all-position or single-position intervention (0/5 both),
reproducing the prior Chronos null. The TSFM's features are predictively
redundant \emph{and} causally quiet at this scale.

\section{Cross-modal synthesis}
\label{sec:synth}

Table~\ref{tab:synth} places the two modalities side by side. Two of the three
legs of the going-in thesis replicate; one dissociates---and the dissociation is
the contribution.

\begin{table}[t]
\centering
\small
\caption{Cross-modal synthesis. Predictive null and deployable selective
predictor replicate in both modalities; the causal contribution is specific to
the language model. \dauroc{} CIs are paired-bootstrap ($B{=}2000$); $p$ is a
one-sided label-permutation test ($B{=}10{,}000$).}
\label{tab:synth}
\begin{tabular}{lccc}
\toprule
 & \multicolumn{2}{c}{\textbf{LM (Pythia)}} & \textbf{TSFM (Chronos)} \\
\cmidrule(lr){2-3}\cmidrule(lr){4-4}
Axis & HellaSwag & SQuAD & ETTh1 \\
\midrule
P1 cheap AUROC                 & 0.509 & 0.590 & 0.654 \\
P2 cheap$+$raw AUROC           & 0.472 & 0.671 & 0.584 \\
P3 cheap$+$\sae{} AUROC        & 0.500 & 0.592 & 0.426 \\
$\Delta$(\sae{}$-$raw)         & $+0.028$ & $-0.079$ & $-0.158$ \\
\quad 95\% CI                  & $[-0.001,.058]$ & $[-0.118,-0.041]$ & $[-0.293,-0.025]$ \\
label-perm $p$ (1-sided)       & $0.033$ & $<10^{-4}$ & $0.011$ \\
causal: all-position           & 5/5 & 5/5 & \textbf{0/5} \\
causal: single-position        & 0/5 & 2/5 & \textbf{0/5} \\
selective: \% oracle AURC      & 2.0\% & 41.3\% & 30.5\% \\
cascade: Pareto-dom.\ points   & 1 & 31 & 5 \\
\bottomrule
\end{tabular}
\end{table}

\paragraph{The one-sentence claim.}
Across an autoregressive LM and an encoder TSFM, TopK-\sae{} features are
predictively redundant for difficulty beyond the strongest cheap baseline
($\Delta$(\sae{}$-$raw)$\le0$, $p\le0.03$) and a cheap-baseline selective
predictor is deployable in both (30--41\% of oracle AURC), \emph{but} the causal
contribution is specific to the LM (5/5 features significant under all-position
patching vs.\ 0/5 on Chronos under either coverage).

\paragraph{Robustness to \sae{} width.}
Re-training each modality at the other's expansion (LM $4\times\!\to\!8\times$,
TSFM $8\times\!\to\!4\times$) and re-probing leaves the null intact:
$\Delta$(\sae{}$-$raw) stays significantly negative on SQuAD ($-0.079\!\to\!-0.058$)
and ETTh1 ($-0.158\!\to\!-0.100$). The \emph{larger} LM \sae{} still does not beat
raw, directly answering ``maybe a bigger \sae{} would have found signal.''

\paragraph{Why the dissociation matters.}
If the TSFM cells had echoed the LM cells, the contribution would be a
cross-modal replication. Because they diverge, the contribution is sharper: the
causal difficulty signal is not a universal property of sparse features but a
property of \emph{this} autoregressive LM, while the \emph{predictive} null and
the \emph{deployable} positive are universal. Prior work did one half each on a
single modality \citep{mishra2026ablation,timesae2026}; the dissociation requires
both halves on both modalities through one pipeline.

\section{Deployable artifact}
\label{sec:deploy}

When the answer to ``do internals help?'' is ``no, but the cheap baseline does,''
the cheap baseline is still deployable. A Platt-recalibrated selective predictor
on the cheap rung captures 41.3\% (LM, raw scale) and 30.5\% (TSFM) of the
oracle's AURC improvement; recalibration on 5-fold out-of-fold predictions cuts
expected calibration error substantially under Platt scaling, which preserves the
ranking (e.g.\ SQuAD P1 ECE $0.366\!\to\!0.123$; ETTh1 P1 $0.284\!\to\!0.156$;
isotonic reaches $0.119$ on ETTh1 but is non-monotone). Routing on the calibrated
score yields Pareto-dominating cost--quality points in both modalities
(31 on SQuAD, 5 on ETTh1), beating always-cheap, always-expensive, and random
routing.

\section{Limitations}
\label{sec:limits}

We probe a single \sae{} family (TopK) at two expansions; a different
dictionary-learning objective could in principle expose signal TopK does not,
though the width sweep argues against a capacity explanation. We use two model
scales per modality and one dataset per LM task and one TSFM dataset; the
near-chance HellaSwag regime limits what can be concluded there. The causal $n$
on the TSFM is bounded by the $167$-window test set; the 0/5 result is a null,
and nulls are weaker than positives, though it reproduces at full sample budget.
Finally, ``difficulty'' is operationalized per modality (binary correctness /
top-quartile CRPS); other operationalizations may behave differently.

\section{Conclusion}
The same sparse features that a probe ranks as ``difficulty-predictive'' add no
predictive power over raw activations in either a language or a time-series
foundation model, yet are causally active in the language model alone. The
deployable lesson is modality-invariant---recalibrate the cheap baseline---while
the scientific lesson is sharp: a causal self-difficulty signal in sparse
features is, at these scales, a property of the autoregressive language model and
not of foundation models in general.

\paragraph{Reproducibility.}
All numbers and figures regenerate from one command
(\texttt{reproduce.sh}); the shared core is unit-tested on synthetic arrays and
contains no model-specific code. Code: \url{https://github.com/nabindev3/fm-difficulty-probe}.

\small
\bibliographystyle{plainnat}
\bibliography{references}

\appendix
\section{Per-modality detail}
The LM adapter (Pythia) and TSFM adapter (Chronos) each own their feature
extractors, hooks, CV scheme, and causal metric; see the per-modality
repositories \url{https://github.com/nabindev3/llm-sae-difficulty} and
\url{https://github.com/nabindev3/tsfm-sae-difficulty} for the full leakage-control
ablations, the 2$\times$2 coverage/fidelity disentanglement on the LM, and the
selective-forecasting calibration curves on the TSFM.

\end{document}
